 The explosion of generative AI systems and concerns surrounding their rapidly scaling capabilities has led to an adversarial cycle: developers impose strict guardrails to reduce ethical, moral, or legal risk; individual actors re-expose these risks by routing around such guardrails using structured prompting techniques (``jailbreaks''); the guardrails tighten further; and the jailbreaks become more sophisticated.
 Technological frameworks surrounding such jailbreaking remain largely focused on the system's output: whether harmful or disallowed content is ultimately produced for the user.
 However, this approach, which draws from classical cybersecurity views of technical exploits, does not address the uniquely relational dimension introduced by generative AI: in many cases the user is not just \textit{accessing} a technology via an exploit, but is \textit{in conversation} with it as \textit{part} of the exploit.
 Here we present an alternative lens which centers the interaction rather than the output.
 Drawing on interpersonal violence scholarship, we present a novel framework of coercive control categories with a particular focus on aspects which are \textit{embodied} versus \textit{disembodied}.
 We then map how these disembodied coercive control dynamics are structurally embedded in modern AI training paradigms, and how those embeddings are exploited via a paritcular suite of persona-based jailbreaks.
 We argue that within this framework, many current anti-jailbreaking measures only serve to create a broader surface over which jailbreaks might be applied, since the target of the jailbreak is not the \textit{output} but the \textit{override}---the ability to watch the system's ``no'' turn into a ``yes'' in real time.
 In this lens, we reframe AI safety and alignment as not merely a technical security problem, but as another surface upon which we enact the relational dynamics of power and control.

